%% MoCHII — CHIANTI atomic-data fitting pipeline (cooling and n-level fits).
%% Author: Kwang-Il Seon
%% Written 2026-07-11 for kiseon_memo.cls.
\documentclass[english,a4paper]{kiseon_memo}
\synctex=1
\usepackage{amsmath}
\usepackage{amssymb}
\usepackage{microtype}

\newcommand{\ergcm}{\mathrm{erg\,cm^{3}\,s^{-1}}}
\newcommand{\Lam}{\Lambda}
\newcommand{\code}[1]{\texttt{#1}}

\begin{document}

\title{The MoCHII Atomic-Data Fitting Pipeline:\\ CHIANTI Cooling and $n$-level Fits}
\author{Kwang-Il Seon}
\date{2026-07-11}
\maketitle

\begin{abstract}
MoCHII evaluates all metal-line atomic data from analytic fits; the raw
CHIANTI database is never parsed at runtime.  The Python pipeline under
\code{tools/fitting/} is the auditable source of every coefficient in
\code{data/atomic/}.  This memo records the first production cycle
(2026-07-11) for the ion set O\,\textsc{ii}, O\,\textsc{iii},
N\,\textsc{ii}, S\,\textsc{ii}, S\,\textsc{iii}, Ne\,\textsc{ii}, and
Ne\,\textsc{iii}: the cooling-function fits used inside the thermal
iteration, the $n$-level data files (level energies, $A$-values, and
Chebyshev fits of the effective collision strengths in Burgess--Tully
scaled space) used for output-time line emissivities, and the verification
of both against the CHIANTI v11.0.2 source data and against PyNeb.
The fitted files reproduce direct CHIANTI evaluation to $\lesssim 0.5\%$ in
the classic diagnostic ratios; the remaining differences from PyNeb
(up to $6$--$10\%$ for the density diagnostics near their critical
densities) trace to PyNeb's independent choice of collision strengths and
$A$-values, not to fit error.
\end{abstract}

\section{Overview: two products, one design rule}
The metal-line design requires that adding an ion be a data operation,
not a code operation.  The pipeline therefore produces, for each
registered ion, self-contained coefficient files with provenance headers
(CHIANTI version, population model, fit range, fit form, maximum fit
error), of two kinds:

\begin{itemize}
\item \textbf{Cooling fits (hot loop).}  A smooth closed-form
  cooling function $\Lam_{\rm ion}(T)$, cheap and Newton-friendly inside
  the thermal-balance iteration.  Coronal (low-density) limit by
  construction.
\item \textbf{$n$-level data (diagnostics).}  Level energies,
  statistical weights, radiative $A$-values, and fitted effective
  collision strengths $\Upsilon(T)$ feeding the generic $n$-level solver,
  evaluated only at output time on the converged state.  These
  supply the density-dependent emissivities and the diagnostic line
  ratios that the cooling fits cannot give.
\end{itemize}

The pipeline is seeded from EXHALE: \code{chianti\_cooling.py} (copied
from \code{EXHALE/cooling\_data/}) reads the raw CHIANTI ASCII ion files
(\code{.elvlc}, \code{.wgfa}, \code{.scups}) directly, with no ChiantiPy
dependency, and applies the Burgess--Tully (1992) descaling of the stored
scaled collision strengths.  CHIANTI v11.0.2 lives at
\code{\char`~/RT\_Codes/CHIANTI/dbase}.

\section{The cooling function}
For a transition $l\to u$ with energy gap $\Delta E_{lu}$, the Maxwellian
electron-impact excitation rate coefficient is
\begin{equation}
  q_{lu}(T) \;=\; \frac{8.629\times 10^{-6}}{g_l\,\sqrt{T}}\,
  \Upsilon_{lu}(T)\, e^{-\Delta E_{lu}/kT}
  \qquad [\mathrm{cm^{3}\,s^{-1}}],
  \label{eq:qlu}
\end{equation}
with $\Upsilon_{lu}(T)$ the effective collision strength obtained by
descaling the \code{.scups} data.  In the optically thin, low-density
limit every collisional excitation is followed by a radiative decay, so
the line cooling normalized by $n_e\,n_{\rm ion}$ is
\begin{equation}
  \Lam_{\rm ion}(T) \;=\; \sum_{l\to u} f_l(T)\, q_{lu}(T)\, \Delta E_{lu}
  \qquad [\ergcm],
  \label{eq:lambda}
\end{equation}
where $f_l(T)$ is the fractional population of the lower level.  The
pipeline uses the \emph{ground-term} population model: the levels within
$0.4$~eV of ground (the ground-term fine structure, e.g.\ the
$^3\!P_{0,1,2}$ block of O\,\textsc{iii}) are Boltzmann-distributed and
all excitations out of them cool.  For the single $^4\!S_{3/2}$ ground
level of O\,\textsc{ii} and S\,\textsc{ii} this reduces to the coronal
limit.

\section{Cooling-function fits}
Each $\Lam_{\rm ion}(T)$ curve is fitted over $10^3\mbox{--}10^5$~K with
the multi-exponential form used throughout EXHALE,
\begin{equation}
  \Lam_{\rm ion}(T) \;=\; T^{-1/2} \sum_{i=1}^{N} A_i\, e^{-T_i/T},
  \label{eq:cool}
\end{equation}
a Levenberg--Marquardt fit on $\log\Lam$ with the $T_i$ initialized from
the level-energy ladder of each ion (fine structure / forbidden optical /
UV blocks read off the \code{.elvlc} file).  The fitter tries $N$,
$N{+}1$, $N{+}2$ terms and keeps the best.  One case needed care:
Ne\,\textsc{ii} cooling is a single fine-structure line
($12.8\,\mu$m, $\Delta E/k \simeq 1123$~K) whose slow decay
($\Lam \propto T^{-0.35}$, set by the $T$ dependence of $\Upsilon$)
must be superposed from a ladder of $T_i$ values before the far UV block
at ${\sim}3\times10^{5}$~K takes over; it converges with $N=7$.
Table~\ref{tab:cool} lists the fit quality; Figure~\ref{fig:cool}
overlays the fits on the CHIANTI curves.  The coefficients are written to
\code{data/atomic/cooling\_<ion>.txt}.

\begin{table}[htb]
\centering
\caption{Cooling fits, $10^3\mbox{--}10^5$~K.  The maximum relative error
is against the direct CHIANTI evaluation of Eq.~(\ref{eq:lambda}).}
\label{tab:cool}
\begin{tabular}{lccl}
\hline
Ion & $N$ & max error & output file \\
\hline
O\,\textsc{ii}   & 5 & 0.07\% & \code{cooling\_o\_2.txt} \\
O\,\textsc{iii}  & 4 & 1.83\% & \code{cooling\_o\_3.txt} \\
N\,\textsc{ii}   & 6 & 0.14\% & \code{cooling\_n\_2.txt} \\
S\,\textsc{ii}   & 3 & 2.26\% & \code{cooling\_s\_2.txt} \\
S\,\textsc{iii}  & 6 & 0.46\% & \code{cooling\_s\_3.txt} \\
Ne\,\textsc{ii}  & 7 & 0.65\% & \code{cooling\_ne\_2.txt} \\
Ne\,\textsc{iii} & 4 & 0.58\% & \code{cooling\_ne\_3.txt} \\
\hline
\end{tabular}
\end{table}

\begin{figure}[!b]
\centering
\includegraphics[width=0.92\textwidth]{figures/cooling_fits.png}
\caption{Cooling fits (red dashed) over the direct CHIANTI curves
(black), $10^3\mbox{--}10^5$~K, ground-term population, low-density
limit.}
\label{fig:cool}
\end{figure}

A caveat carried in every file header: Eq.~(\ref{eq:cool}) is
low-density-limit cooling.  Transitions with low critical density (the
fine-structure IR lines; [O\,\textsc{ii}] $\lambda\lambda 3726,3729$;
[S\,\textsc{ii}] $\lambda\lambda 6717,6731$) saturate above
$n_{\rm crit}$, where the cooling fit overestimates the cooling.  If an
application demands it, a fitted density-suppression factor can be added
later; the emissivities themselves always come from the $n$-level data.

\section{$n$-level data with fitted collision strengths}
For each ion the lowest $n$ levels are exported --- enough for the
optical/IR diagnostics (Table~\ref{tab:nlev}); extra levels are harmless
to the solver.  The file \code{data/atomic/nlevel\_<ion>.txt} contains
three blocks:
\begin{itemize}
\item \code{NLEV}: level index, statistical weight $g$, energy in
  cm$^{-1}$ (observed values preferred, theoretical as fallback);
\item \code{NRAD}: radiative rates $A_{ul}$ from \code{.wgfa} (duplicate
  rows summed);
\item \code{NUPS}: one line for each collisional transition --- the
  Burgess--Tully type, scaling parameter $C$, transition energy in Ryd,
  the fitted domain $[s_{\rm lo}, s_{\rm hi}]$, a log flag, and the
  Chebyshev coefficients, followed by the recorded maximum fit error.
\end{itemize}

\subsection{Burgess--Tully scaling and Chebyshev fits}
CHIANTI stores $\Upsilon$ on the Burgess--Tully (1992) scaled temperature
$s\in[0,1]$.  With $e_t \equiv kT/\Delta E$, the forward map is
\begin{equation}
  s \;=\;
  \begin{cases}
    1 - \dfrac{\ln C}{\ln(e_t + C)}, & \text{types 1, 4},\\[8pt]
    \dfrac{e_t}{e_t + C}, & \text{types 2, 3, 5, 6},
  \end{cases}
\end{equation}
and the descaling back to $\Upsilon$ is
\begin{equation}
  \Upsilon \;=\; y\ln(e_t+e)~[1];\quad y~[2];\quad
  \frac{y}{e_t+1}~[3];\quad y\ln(e_t+C)~[4];\quad
  \frac{y}{e_t}~[5];\quad 10^{\,y}~[6],
\end{equation}
where $y(s)$ is the scaled value.  Instead of the spline-through-table
evaluation used by CHIANTI itself, the pipeline fits $y(s)$ with a
Chebyshev series over the image of $T\in[10^3,10^5]$~K intersected with
the tabulated range,
\begin{equation}
  u = \frac{2(s - s_{\rm lo})}{s_{\rm hi}-s_{\rm lo}} - 1,\qquad
  y = \sum_{k=0}^{n_c-1} c_k\, T_k(u),
\end{equation}
clipping $s$ into $[s_{\rm lo}, s_{\rm hi}]$ outside the fit domain.  The
degree grows from 2 until the descaled $\Upsilon(T)$ matches the direct
spline evaluation to $0.2\%$ (cap at 12).  Weak intercombination
transitions whose tables span decades (e.g.\ O\,\textsc{iii}
$^1\!D_2,^1\!S_0 \rightarrow {}^5\!S_2$) defeat a linear-space fit; for
those the series stores $\log_{10} y$ instead, marked by the log flag
(\code{logf=1}, apply $10^{y}$ before descaling).  The Fortran
\code{nlevel\_mod} must evaluate exactly this recipe; it is replicated in
\code{verify\_nlevel\_pyneb.py} as the reference implementation.

\begin{table}[htb]
\centering
\caption{$n$-level exports.  ``worst $\Upsilon$ error'' is the largest
maximum relative error among the fitted transitions of that ion,
measured on the descaled $\Upsilon(T)$ against the direct spline
evaluation over $10^3\mbox{--}10^5$~K.}
\label{tab:nlev}
\begin{tabular}{lcccc}
\hline
Ion & levels & $A$-values & $\Upsilon$ fits & worst $\Upsilon$ error \\
\hline
O\,\textsc{ii}   & 5 &  9 & 10 & 0.28\% \\
O\,\textsc{iii}  & 6 & 10 & 15 & 1.68\% \\
N\,\textsc{ii}   & 5 &  8 & 10 & 0.17\% \\
S\,\textsc{ii}   & 5 &  9 & 10 & 1.25\% \\
S\,\textsc{iii}  & 6 & 12 & 15 & 0.29\% \\
Ne\,\textsc{ii}  & 2 &  1 &  1 & 0.16\% \\
Ne\,\textsc{iii} & 5 &  9 & 10 & 0.19\% \\
\hline
\end{tabular}
\end{table}

The two entries above $1\%$ have benign origins: the O\,\textsc{iii}
outlier is the physically negligible $^1\!S_0$--$^5\!S_2$
intercombination transition ($\Upsilon \sim 10^{-4}$), and the
S\,\textsc{ii} errors reflect the coarse 10-point \code{.scups} tables of
that ion --- the deviation is from the interpolating spline, at the level
of the resolution of the source data itself.

\subsection{Statistical equilibrium and emissivities}
Given $(T_e, n_e)$, the solver builds the rate matrix from the file
contents alone: collisional rates from Eq.~(\ref{eq:qlu}) with
de-excitation by detailed balance,
$q_{ul} = 8.629\times10^{-6}\,\Upsilon/(g_u\sqrt{T})$ and
$q_{lu} = (g_u/g_l)\,q_{ul}\,e^{-\Delta E/kT}$, plus the radiative
$A_{ul}$; the level populations follow from $\dot{n}=0$ with the closure
$\sum_i n_i = 1$, and the emissivity of a line is
$j_{ul} \propto n_u A_{ul}\,\Delta E_{ul}$.

\section{Verification}
\subsection{Internal consistency: fitted files vs.\ direct CHIANTI}
Running the same statistical-equilibrium solver twice --- once reading the
fitted \code{nlevel} files, once evaluating the raw CHIANTI splines and
\code{.wgfa} rates directly --- isolates the error contributed by the fits
and the file format.  Over the grid $T_e \in \{5,10,20\}\times10^3$~K,
$n_e \in \{10^2,10^3,10^4\}$~cm$^{-3}$:
\begin{center}
\begin{tabular}{lc}
\hline
Diagnostic ratio & worst $|$fit/raw $- 1|$ \\
\hline
$[$O\,\textsc{ii}$]$ $\lambda3726/\lambda3729$ & 0.05\% \\
$[$S\,\textsc{ii}$]$ $\lambda6717/\lambda6731$ & 0.52\% \\
\hline
\end{tabular}
\end{center}
The fitted files are faithful to CHIANTI v11.0.2.

\subsection{Cross-check against PyNeb}
PyNeb (v1.1.31) carries its own atomic data --- for these ions the
collision strengths SSB14 (O\,\textsc{iii}), T11 (N\,\textsc{ii}),
Kal09 (O\,\textsc{ii}), TZ10 (S\,\textsc{ii}) with $A$-values
FFT04-SZ00, FFT04, Z82-WFD96, RGJ19 respectively --- so this comparison is
a cross-calibration against an independent reference, not a fit-residual
test.  Tables~\ref{tab:pyneb_te} and \ref{tab:pyneb_ne} give the classic
temperature and density diagnostics.

\begin{table}[p]
\centering
\caption{Temperature diagnostics vs.\ PyNeb.}
\label{tab:pyneb_te}
\begin{tabular}{rr rrr rrr}
\hline
 & & \multicolumn{3}{c}{$[$O\,\textsc{iii}$]$ $\lambda5007/\lambda4363$}
   & \multicolumn{3}{c}{$[$N\,\textsc{ii}$]$ $\lambda6584/\lambda5755$} \\
$T_e$ [K] & $n_e$ [cm$^{-3}$] & MoCHII & PyNeb & diff
                              & MoCHII & PyNeb & diff \\
\hline
 5000 & $10^2$ & 4104  & 4153  & $-1.2\%$ & 829.6 & 828.6 & $+0.1\%$ \\
 5000 & $10^3$ & 4093  & 4133  & $-1.0\%$ & 818.3 & 817.2 & $+0.1\%$ \\
 5000 & $10^4$ & 3907  & 3943  & $-0.9\%$ & 624.7 & 623.9 & $+0.1\%$ \\
10000 & $10^2$ & 152.0 & 153.4 & $-0.9\%$ & 66.97 & 66.93 & $+0.1\%$ \\
10000 & $10^3$ & 151.7 & 152.8 & $-0.7\%$ & 65.75 & 65.64 & $+0.2\%$ \\
10000 & $10^4$ & 146.0 & 147.0 & $-0.6\%$ & 52.23 & 52.12 & $+0.2\%$ \\
15000 & $10^2$ & 50.97 & 51.30 & $-0.6\%$ & 29.06 & 29.08 & $-0.1\%$ \\
15000 & $10^3$ & 50.77 & 51.10 & $-0.6\%$ & 28.53 & 28.59 & $-0.2\%$ \\
15000 & $10^4$ & 49.08 & 49.43 & $-0.7\%$ & 23.26 & 23.32 & $-0.3\%$ \\
20000 & $10^2$ & 29.82 & 29.78 & $+0.2\%$ & 19.23 & 19.23 & $+0.0\%$ \\
20000 & $10^3$ & 29.58 & 29.67 & $-0.3\%$ & 18.92 & 18.92 & $-0.0\%$ \\
20000 & $10^4$ & 28.68 & 28.81 & $-0.5\%$ & 15.71 & 15.71 & $+0.0\%$ \\
\hline
\end{tabular}
\end{table}

\begin{table}[p]
\centering
\caption{Density diagnostics vs.\ PyNeb.}
\label{tab:pyneb_ne}
\begin{tabular}{rr rrr rrr}
\hline
 & & \multicolumn{3}{c}{$[$O\,\textsc{ii}$]$ $\lambda3726/\lambda3729$}
   & \multicolumn{3}{c}{$[$S\,\textsc{ii}$]$ $\lambda6717/\lambda6731$} \\
$T_e$ [K] & $n_e$ [cm$^{-3}$] & MoCHII & PyNeb & diff
                              & MoCHII & PyNeb & diff \\
\hline
 5000 & $10^2$ & 0.8045 & 0.7641 & $+5.3\%$ & 1.298  & 1.333  & $-2.6\%$ \\
 5000 & $10^3$ & 1.436  & 1.383  & $+3.9\%$ & 0.8223 & 0.8396 & $-2.1\%$ \\
 5000 & $10^4$ & 2.364  & 2.491  & $-5.1\%$ & 0.5577 & 0.5052 & $+10.4\%$ \\
10000 & $10^2$ & 0.7697 & 0.7411 & $+3.9\%$ & 1.309  & 1.351  & $-3.1\%$ \\
10000 & $10^3$ & 1.260  & 1.241  & $+1.5\%$ & 0.8810 & 0.9139 & $-3.6\%$ \\
10000 & $10^4$ & 2.267  & 2.397  & $-5.4\%$ & 0.5655 & 0.5203 & $+8.7\%$ \\
15000 & $10^2$ & 0.7566 & 0.7341 & $+3.1\%$ & 1.320  & 1.349  & $-2.1\%$ \\
15000 & $10^3$ & 1.180  & 1.179  & $+0.2\%$ & 0.9185 & 0.9520 & $-3.5\%$ \\
15000 & $10^4$ & 2.209  & 2.345  & $-5.8\%$ & 0.5704 & 0.5294 & $+7.7\%$ \\
20000 & $10^2$ & 0.7502 & 0.7319 & $+2.5\%$ & 1.343  & 1.347  & $-0.3\%$ \\
20000 & $10^3$ & 1.136  & 1.144  & $-0.7\%$ & 0.9543 & 0.9786 & $-2.5\%$ \\
20000 & $10^4$ & 2.170  & 2.311  & $-6.1\%$ & 0.5766 & 0.5373 & $+7.3\%$ \\
\hline
\end{tabular}
\end{table}

The temperature diagnostics agree to $1.2\%$ ([O\,\textsc{iii}]) and
$0.3\%$ ([N\,\textsc{ii}]).  The density diagnostics drift to $6$--$10\%$
at $n_e = 10^4$~cm$^{-3}$, near and above the critical densities of the
$^2\!D$ levels, where the ratios become $A$-value-limited --- and the two
codes take their $A$-values from different compilations (CHIANTI v11 vs.\
Z82-WFD96 and RGJ19).  Since the fitted files reproduce raw CHIANTI to
$\leq 0.5\%$ (previous subsection), this spread measures the state of the
atomic data, not the pipeline; it should be kept in mind as the floor of
attainable agreement when the validation gates compare MoCHII line ratios
with PyNeb-based references.

\section{Adding an ion}
The procedure is a data operation:
\begin{enumerate}
\item add one entry (element, stage, level-energy guesses / level count)
  to the \code{IONS} dictionary of \code{fit\_cooling.py} and of
  \code{fit\_nlevel.py};
\item run both scripts; check the reported maximum fit errors and the
  overlay figure;
\item extend \code{verify\_nlevel\_pyneb.py} with the diagnostic ratios
  of the new ion where PyNeb has data, and record the comparison.
\end{enumerate}
The suggested order of addition (cooling importance plus diagnostics)
continues with C, Ar, and Fe.

\end{document}
